\clearpage
\section*{Problem 1: Orthogonal Matching Pursuit (OMP)}
\addcontentsline{toc}{section}{Problem 1: Orthogonal Matching Pursuit}
\markboth{Problem 1}{Problem 1}

\textbf{Problem Statement:} 

\subsection*{Part 0: Preliminary Information and Data}

In part 2, use the matrix
$$A = \begin{bmatrix} 
1 & -1/2 & -1/2 \\
0 & \sqrt{3}/2 & -\sqrt{3}/2 \\
1 & 3 & 3
\end{bmatrix}$$

In part 3, use the matrix $B$ that is the $2 \times 3$ matrix consisting of the first two rows of $A$.

\subsection*{Singular Value Decomposition (SVD) and Pseudo-Inverse}

Given a matrix $M_{m \times n}$ with $\text{rank}(M) = r$, write
$$M = U\Sigma V^T$$
where $U_{m \times m}$ and $V_{n \times n}$ are orthogonal matrices, and $\Sigma_{m \times n}$ is a diagonal matrix with nonzero diagonal entries $\sigma_1 \geq \sigma_2 \geq \ldots \geq \sigma_r > 0$ called singular values of $M$.

\textbf{Definition.} Given $M$ as above, the pseudo-inverse $M^+$ is defined by
$$M^+ = V\Sigma^+ U^T$$
where $\Sigma^+_{n \times m}$ is the diagonal matrix whose non-zero entries are $\frac{1}{\sigma_1}, \frac{1}{\sigma_2}, \ldots, \frac{1}{\sigma_r}$.

\subsection*{Thresholding Operator}

The \textbf{hard thresholding operator} $H_s$ is defined as follows:
$$\mathbf{x} \to H_s(\mathbf{x}) \in \mathbb{R}^n$$

$H_s(\mathbf{x})$ sets all entries of $\mathbf{x}$ to zero except for $s$ largest (in absolute value) entries.

\subsection*{The Matrix $A_\mathcal{S}$}

Given a matrix $A$ and an index set $\mathcal{S}$, denote by $A_\mathcal{S}$ the matrix consisting of columns of $A$ with index in $\mathcal{S}$.

\noindent\rule{\textwidth}{0.4pt}

\vspace{1em}

\subsection*{Assignments and Questions}

\begin{enumerate}
\item Implement the pseudocode for the Orthogonal Matching Pursuit (OMP). The pseudocode is contained in a separate handwritten note attached to the Final project assignment. You may use the standard Matlab (or other language) commands to compute SVD and pseudo-inverse.

\item Run the code for $\mathbf{b} = A\mathbf{e}_1$ with $\mathbf{e}_1 = (1, 0, 0)^T$. See how well the obtained vector $\mathbf{x}$ matches $\mathbf{e}_1$.

\item Run the code again, this time using the matrix $B$ instead of $A$.

\item Compare the results obtained in parts 2, 3. Explain the connection with compressed sensing and sparse recovery.
\end{enumerate}

\vspace{1em}

\textbf{Solution:}

% TODO: Add solution here

\vspace{2em}
\noindent\rule{\textwidth}{0.4pt}
